\documentclass[12pt,a4paper]{article}

\usepackage[margin=1in]{geometry}
\usepackage{amsmath, amssymb, amsthm}
\usepackage{enumitem}
\usepackage{booktabs}

\setlength{\parindent}{0pt}
\setlength{\parskip}{6pt}

\newtheoremstyle{notestyle}
  {8pt}
  {8pt}
  {\normalfont}
  {}
  {\bfseries}
  {.}
  {0.5em}
  {}

\theoremstyle{notestyle}
\newtheorem{example}{Example}

\newenvironment{answer}
{\par\textbf{Answer.}\ }
{\hfill$\square$\par}

\title{\textbf{Finding the Limiting Distribution Using Eigenvalues}}
\author{MA4034: Time Series and Stochastic Processes}
\date{}

\begin{document}
\maketitle

\section{Goal}

For the ON/OFF Markov chain,
\[
P=
\begin{pmatrix}
1-p & p\\
q & 1-q
\end{pmatrix},
\]
we want to find
\[
\lim_{n\to\infty}P^n.
\]

The states are ordered as
\[
(0,1)=(\text{OFF},\text{ON}).
\]

\section{Step 1: Find the Eigenvalues}

The eigenvalues of $P$ are
\[
\lambda_1=1,
\qquad
\lambda_2=1-p-q.
\]

So the diagonal matrix is
\[
D=
\begin{pmatrix}
1 & 0\\
0 & 1-p-q
\end{pmatrix}.
\]

\section{Step 2: Find the Eigenvectors}

To find eigenvectors, solve
\[
(P-\lambda I)v=0.
\]

Let
\[
v=
\begin{pmatrix}
x\\
y
\end{pmatrix}.
\]

\subsection{Eigenvector for $\lambda_1=1$}

Solve
\[
(P-I)v=0.
\]

Now
\[
P-I
=
\begin{pmatrix}
1-p & p\\
q & 1-q
\end{pmatrix}
-
\begin{pmatrix}
1 & 0\\
0 & 1
\end{pmatrix}.
\]

Therefore,
\[
P-I
=
\begin{pmatrix}
-p & p\\
q & -q
\end{pmatrix}.
\]

Now solve
\[
\begin{pmatrix}
-p & p\\
q & -q
\end{pmatrix}
\begin{pmatrix}
x\\
y
\end{pmatrix}
=
\begin{pmatrix}
0\\
0
\end{pmatrix}.
\]

The first row gives
\[
-px+py=0.
\]

So
\[
p(y-x)=0.
\]

Assuming $p\neq0$,
\[
y=x.
\]

Choose
\[
x=1.
\]

Then
\[
y=1.
\]

So one eigenvector is
\[
v_1=
\begin{pmatrix}
1\\
1
\end{pmatrix}.
\]

\subsection{Eigenvector for $\lambda_2=1-p-q$}

Solve
\[
(P-\lambda_2 I)v=0.
\]

That is,
\[
\left(P-(1-p-q)I\right)v=0.
\]

Now
\[
P-(1-p-q)I
=
\begin{pmatrix}
1-p & p\\
q & 1-q
\end{pmatrix}
-
\begin{pmatrix}
1-p-q & 0\\
0 & 1-p-q
\end{pmatrix}.
\]

Therefore,
\[
P-(1-p-q)I
=
\begin{pmatrix}
q & p\\
q & p
\end{pmatrix}.
\]

Now solve
\[
\begin{pmatrix}
q & p\\
q & p
\end{pmatrix}
\begin{pmatrix}
x\\
y
\end{pmatrix}
=
\begin{pmatrix}
0\\
0
\end{pmatrix}.
\]

This gives
\[
qx+py=0.
\]

So
\[
py=-qx.
\]

Hence
\[
y=-\frac{q}{p}x.
\]

Choose
\[
x=p.
\]

Then
\[
y=-q.
\]

So one eigenvector is
\[
v_2=
\begin{pmatrix}
p\\
-q
\end{pmatrix}.
\]

\section{Step 3: Form the Matrix $R$}

The matrix $R$ is formed by putting the eigenvectors as columns:
\[
R=
\begin{pmatrix}
| & |\\
v_1 & v_2\\
| & |
\end{pmatrix}.
\]

So
\[
R=
\begin{pmatrix}
1 & p\\
1 & -q
\end{pmatrix}.
\]

Then
\[
P=RDR^{-1}.
\]

\section{Step 4: Find $R^{-1}$}

For a $2\times2$ matrix
\[
\begin{pmatrix}
a & b\\
c & d
\end{pmatrix},
\]
the inverse is
\[
\frac{1}{ad-bc}
\begin{pmatrix}
d & -b\\
-c & a
\end{pmatrix}.
\]

Here
\[
R=
\begin{pmatrix}
1 & p\\
1 & -q
\end{pmatrix}.
\]

So
\[
a=1,\qquad b=p,\qquad c=1,\qquad d=-q.
\]

The determinant is
\[
ad-bc=(1)(-q)-(p)(1).
\]

Thus,
\[
ad-bc=-q-p=-(p+q).
\]

Therefore,
\[
R^{-1}
=
\frac{1}{-(p+q)}
\begin{pmatrix}
-q & -p\\
-1 & 1
\end{pmatrix}.
\]

Multiplying the negative sign inside,
\[
R^{-1}
=
\frac{1}{p+q}
\begin{pmatrix}
q & p\\
1 & -1
\end{pmatrix}.
\]

\section{Step 5: Find $P^n$}

Since
\[
P=RDR^{-1},
\]
we have
\[
P^n=RD^nR^{-1}.
\]

Now
\[
D=
\begin{pmatrix}
1 & 0\\
0 & 1-p-q
\end{pmatrix}.
\]

Therefore,
\[
D^n=
\begin{pmatrix}
1^n & 0\\
0 & (1-p-q)^n
\end{pmatrix}.
\]

Since
\[
1^n=1,
\]
we get
\[
D^n=
\begin{pmatrix}
1 & 0\\
0 & (1-p-q)^n
\end{pmatrix}.
\]

Thus,
\[
P^n=
\begin{pmatrix}
1 & p\\
1 & -q
\end{pmatrix}
\begin{pmatrix}
1 & 0\\
0 & (1-p-q)^n
\end{pmatrix}
\frac{1}{p+q}
\begin{pmatrix}
q & p\\
1 & -1
\end{pmatrix}.
\]

\section{Step 6: Multiply to Get the $n$-Step Matrix}

First multiply
\[
R D^n.
\]

\[
\begin{pmatrix}
1 & p\\
1 & -q
\end{pmatrix}
\begin{pmatrix}
1 & 0\\
0 & (1-p-q)^n
\end{pmatrix}
=
\begin{pmatrix}
1 & p(1-p-q)^n\\
1 & -q(1-p-q)^n
\end{pmatrix}.
\]

Now multiply by $R^{-1}$:
\[
P^n
=
\frac{1}{p+q}
\begin{pmatrix}
1 & p(1-p-q)^n\\
1 & -q(1-p-q)^n
\end{pmatrix}
\begin{pmatrix}
q & p\\
1 & -1
\end{pmatrix}.
\]

Compute each entry.

First row, first column:
\[
q+p(1-p-q)^n.
\]

First row, second column:
\[
p-p(1-p-q)^n.
\]

Second row, first column:
\[
q-q(1-p-q)^n.
\]

Second row, second column:
\[
p+q(1-p-q)^n.
\]

Therefore,
\[
P^n
=
\frac{1}{p+q}
\begin{pmatrix}
q+p(1-p-q)^n & p-p(1-p-q)^n\\
q-q(1-p-q)^n & p+q(1-p-q)^n
\end{pmatrix}.
\]

This can also be written as
\[
P^n
=
\frac{1}{p+q}
\begin{pmatrix}
q & p\\
q & p
\end{pmatrix}
+
\frac{(1-p-q)^n}{p+q}
\begin{pmatrix}
p & -p\\
-q & q
\end{pmatrix}.
\]

\section{Step 7: Take the Limit}

Assume
\[
|1-p-q|<1.
\]

Then
\[
\lim_{n\to\infty}(1-p-q)^n=0.
\]

So in
\[
P^n
=
\frac{1}{p+q}
\begin{pmatrix}
q & p\\
q & p
\end{pmatrix}
+
\frac{(1-p-q)^n}{p+q}
\begin{pmatrix}
p & -p\\
-q & q
\end{pmatrix},
\]
the second part goes to zero.

Therefore,
\[
\lim_{n\to\infty}P^n
=
\frac{1}{p+q}
\begin{pmatrix}
q & p\\
q & p
\end{pmatrix}.
\]

So
\[
\boxed{
\lim_{n\to\infty}P^n
=
\begin{pmatrix}
\frac{q}{p+q} & \frac{p}{p+q}\\
\frac{q}{p+q} & \frac{p}{p+q}
\end{pmatrix}.
}
\]

\section{Step 8: Read the Limiting Distribution}

Both rows of the limiting matrix are the same:
\[
\left(\frac{q}{p+q},\frac{p}{p+q}\right).
\]

Therefore the limiting distribution is
\[
\boxed{
\pi=
\left(\frac{q}{p+q},\frac{p}{p+q}\right).
}
\]

Using state order
\[
(\text{OFF},\text{ON}),
\]
this means
\[
\Prob(\text{OFF in long run})=\frac{q}{p+q},
\]
and
\[
\Prob(\text{ON in long run})=\frac{p}{p+q}.
\]

\section{Short Summary}

\[
\begin{array}{c|c}
\text{Step} & \text{What to do}\\
\hline
1 & \text{Find eigenvalues}\\
2 & \text{Find eigenvectors}\\
3 & \text{Put eigenvectors into }R\\
4 & \text{Find }R^{-1}\\
5 & \text{Use }P^n=RD^nR^{-1}\\
6 & \text{Take }n\to\infty\\
7 & \text{Read the repeated row as }\pi
\end{array}
\]

\section{Practice Question}

\begin{example}
Let
\[
p=\frac34,
\qquad
q=\frac12.
\]
Find the limiting distribution of the ON/OFF chain.
\end{example}

\begin{answer}
The limiting distribution is
\[
\pi=
\left(\frac{q}{p+q},\frac{p}{p+q}\right).
\]

Substitute
\[
p=\frac34,
\qquad
q=\frac12.
\]

Then
\[
p+q=\frac34+\frac12=\frac34+\frac24=\frac54.
\]

So
\[
\pi_0=\frac{q}{p+q}
=
\frac{\frac12}{\frac54}
=
\frac12\cdot\frac45
=
\frac25.
\]

And
\[
\pi_1=\frac{p}{p+q}
=
\frac{\frac34}{\frac54}
=
\frac34\cdot\frac45
=
\frac35.
\]

Therefore,
\[
\boxed{
\pi=\left(\frac25,\frac35\right).
}
\]

This means the long-run probability of OFF is $2/5$, and the long-run probability of ON is $3/5$.
\end{answer}

\end{document}
